\section{Information sharing between coordinates}
\label{sec:ntm}

In previous sections we considered a coordinatewise architecture, which
corresponds by analogy to a learned version of RMSprop or ADAM.  Although
diagonal methods are quite effective in practice, we can also consider learning
more sophisticated optimizers that take the correlations between coordinates
into effect. To this end, we introduce a mechanism allowing different LSTMs to
communicate with each other.

\subsection*{Global averaging cells}

The simplest solution is to designate a subset of the cells in each LSTM layer
for communication. These cells operate like normal LSTM cells, but their
outgoing activations are averaged at each step across all coordinates.  These
\emph{global averaging cells} (GACs) are sufficient to allow the networks to
implement L2 gradient clipping \citep{bengio:2013} assuming that each LSTM can
compute the square of the gradient. This architecture is denoted as an LSTM+GAC
optimizer.

\subsection*{NTM-BFGS optimizer}

We also consider augmenting the LSTM+GAC architecture with an external memory
that is shared between coordinates. Such a memory, if appropriately designed
could allow the optimizer to learn algorithms similar to (low-memory)
approximations to Newton's method, e.g.\ (L-)BFGS \citep[see][]{nocedal:2006}.
The reason for this interpretation is that such methods can be seen as a set of
independent processes working coordinatewise, but communicating through the
inverse Hessian approximation stored in the memory.  We designed a memory
architecture that, in theory, allows the network to simulate (L-)BFGS,
motivated by the approximate Newton method BFGS, named for Broyden, Fletcher,
Goldfarb, and Shanno.  We call this architecture an \emph{NTM-BFGS optimizer},
because its use of external memory is similar to the Neural Turing Machine
\citep{graves:2014}.  The pivotal differences between our construction and the
NTM are (1) our memory allows only low-rank updates; (2) the controller
(including read/write heads) operates coordinatewise.

%The construction can be describec by deconstructing an inverse Hessian BFGS algorithm and inserting learnable components.
%in a manner to how we constructed the diagonal variant.
In BFGS an explicit estimate of the full (inverse) Hessian is built up from the sequence of observed gradients.
We can write a skeletonized version of the BFGS algorithm, using $M_t$ to represent the inverse Hessian approximation at iteration $t$, as follows
\begin{align*}
g_t &= \operatorname{read}(M_t, \theta_t) \\
\theta_{t+1} &= \theta_{t} + g_t \\
M_{t+1} &= \operatorname{write}(M_t, \theta_t, g_t)
\,.
\end{align*}
Here we have packed up all of the details of the BFGS algorithm into the
suggestively named $\operatorname{read}$ and $\operatorname{write}$ operations,
which operate on the inverse Hessian approximation $M_t$.  In BFGS these
operations have specific forms, for example $\operatorname{read}(M_t, \theta_t)
= -M_t\nabla f(\theta_t)$ is a specific matrix-vector multiplication and
the BFGS write operation corresponds to a particular low-rank update of $M_t$.

In this work we preserve the structure of the BFGS updates, but discard their
particular form. More specifically the $\operatorname{read}$ operation remains
a matrix-vector multiplication but the form of the vector used is learned.
Similarly, the $\operatorname{write}$ operation remains a low-rank update, but the
vectors involved are also learned. Conveniently, this structure of interaction with
a large dynamically updated state corresponds in a fairly direct way to the
architecture of a Neural Turing Machine (NTM), where $M_t$ corresponds to the
NTM memory \citep{graves:2014}.

Our NTM-BFGS optimizer uses an LSTM+GAC as a controller;
%with the same structure as was used in
%the \piLSTM{} optimizer, with a single small LSTM operating on each coordinate of
%the model
however, instead of producing the update directly we attach one or
more read and write heads to the controller.
Each read head produces a read
vector $r_{t}$ which is combined with the memory to produce a read result $i_t$ which is fed back into the controller at the following time
step.  Each write head produces two outputs, a left write vector $a_t$
and a right write vector $b_t$. The two write vectors are used to update the
memory state by accumulating their outer product. The read and write operation for a single head is
diagrammed in Figure~\ref{fig:ntm-bfgs} and the way read and write heads are attached to the controller
is depicted in Figure~\ref{fig:ntm-controller}.

In can be shown that NTM-BFGS with one read head and 3 write heads can simulate inverse Hessian BFGS assuming that the controller
can compute arbitrary (coordinatewise) functions and have access to 2 GACs.

\begin{figure}
\centering
\includegraphics[width=0.54\linewidth]{figures/ntm-read.pdf}
\qquad
\includegraphics[width=0.38\linewidth]{figures/ntm-write.pdf}
\caption{\textbf{Left:} NTM-BFGS read operation. \textbf{Right:} NTM-BFGS write operation.}
\label{fig:ntm-bfgs}
\end{figure}

\begin{figure}
\centering
\includegraphics[width=0.6\linewidth,valign=t]{figures/ntm-controller.pdf}
\qquad
\includegraphics[width=0.32\linewidth,valign=t]{figures/ntm-controller-element.pdf}
\caption{\textbf{Left:} Interaction between the controller and the external
memory in NTM-BFGS. The controller is composed of replicated coordinatewise
LSTMs (possibly with GACs), but the read and write operations are global across all
coordinates.
\textbf{Right:}  A single LSTM for the $k$th coordinate in the NTM-BFGS
controller. Note that here we have dropped the time index $t$ to simplify
notation.}
\label{fig:ntm-controller}
\end{figure}

\subsection*{NTM-L-BFGS optimizer}

In cases where memory is constrained we can follow the example of
L-BFGS and maintain a low rank approximation of the
full memory (\emph{vis}.\ inverse Hessian).
The simplest way to do this is to store a sliding history of the left and right write vectors,
allowing us to form the matrix vector multiplication required by the read operation efficiently.




% We have also considered replacing
% \begin{align*}
% \nabla^k &\rightarrow (\tanh(\nabla^k e^{s_1})\,, \tanh(\nabla^k e^{s_2})\,,\ldots)
% \end{align*}
% for a number of different scales $s_j$. Both of these transformations have the
% effect of keeping the gradients to the optimizer network on a similar scale
% throughout the optimization process.  In the NTM-BFGS variant we pre-process
% all controller inputs in this way, including the result of the read operation.

% \paragraph{Cost shaping}  Changing the optimizer objective so that
% \begin{align*}
% h(\theta_t) \rightarrow -\exp(-w_t h(\theta_t))
% \end{align*}
% helps to shift emphasis of the optimizer away from the early stages of optimization, allowing it to focus more on later stages.

% \paragraph{Learned trust region} Using an L2 penalty on the magnitude of the updates proposed by the optimizer helps bias it towards producing smaller updates, and helps guard against a single bad update causing a catastrophic failure.

